\chapter{Resultados y Discusión}

Este capítulo presenta los resultados obtenidos de los tres módulos de inteligencia artificial
implementados en el backend del chatbot terapéutico \textit{Seren}: el detector de lenguaje
depresivo (\texttt{DepressionDetectorService}), el analizador narrativo PHQ-9
(\texttt{PHQ9Service}) y el sistema de evaluación conversacional PHQ-9
(\texttt{ConversationalPHQ9Service}). Todos los módulos utilizan el modelo de lenguaje
LLaMA 3.1 8B a través de Ollama como motor de inferencia, con prompt engineering
específico para el dominio de salud mental en español.

% ─────────────────────────────────────────────────────────────────────────────
\section{Desempeño de los Módulos de IA}
% ─────────────────────────────────────────────────────────────────────────────

\subsection{Módulo de Detección de Lenguaje Depresivo}

El detector analiza cada mensaje del usuario en tiempo real (background task) y produce
cuatro salidas: un indicador binario de lenguaje depresivo, un nivel de confianza (0–100),
un nivel de riesgo (bajo/medio/alto/severo) y palabras clave detectadas.

\begin{table}[H]
\centering
\caption{Desempeño interno del detector de lenguaje depresivo -- conjunto de prueba (n=300)}
\label{tab:depression_internal}
\begin{tabular}{lcccc}
\toprule
\textbf{Métrica} & \textbf{Valor} & \textbf{IC 95\%} & \textbf{Umbral} & \textbf{Evaluación} \\
\midrule
Accuracy         & 88.33\%        & [84.1, 91.9]    & --              & Buena     \\
Sensitivity      & \textbf{93.2\%}& [88.6, 96.3]    & 0.50            & Excelente \\
Specificity      & 82.4\%         & [76.5, 87.3]    & 0.50            & Buena     \\
F1-Score         & 0.891          & [0.855, 0.923]  & --              & Muy buena \\
AUC-ROC          & \textbf{0.941} & [0.912, 0.966]  & --              & Excelente \\
\bottomrule
\end{tabular}
\end{table}

\textbf{Observaciones clave:}
\begin{itemize}
    \item La sensitivity elevada (93.2\%) refleja la prioridad clínica de no perder casos
          depresivos.
    \item La especificidad más baja (82.4\%) es aceptable dado que los falsos positivos solo
          desencadenan la evaluación PHQ-9 conversacional, no un diagnóstico.
    \item El AUC de 0.941 indica excelente capacidad discriminativa del modelo LLaMA 3.1 8B
          con el prompt diseñado.
\end{itemize}

\subsection{Módulo de Análisis Narrativo PHQ-9}

El módulo analiza textos largos y predice, para cada uno de los 9 síntomas PHQ-9, si
está presente con un umbral de confianza mínimo de 60\% (\texttt{MIN\_CONFIDENCE = 0.60}).

\begin{table}[H]
\centering
\caption{Desempeño por síntoma PHQ-9 -- análisis narrativo (n=150 narrativas)}
\label{tab:phq9_symptom}
\begin{tabular}{llccc}
\toprule
\textbf{Q} & \textbf{Síntoma}             & \textbf{Precision} & \textbf{Recall} & \textbf{F1} \\
\midrule
Q1 & Poco interés/placer          & 0.891 & 0.923 & 0.907 \\
Q2 & Sentirse deprimido           & 0.912 & 0.947 & 0.929 \\
Q3 & Problemas de sueño           & 0.878 & 0.891 & 0.884 \\
Q4 & Fatiga / poca energía        & 0.903 & 0.917 & 0.910 \\
Q5 & Cambios en apetito           & 0.856 & 0.872 & 0.864 \\
Q6 & Sentirse fracasado           & 0.834 & 0.865 & 0.849 \\
Q7 & Problemas concentración      & 0.871 & 0.889 & 0.880 \\
Q8 & Agitación psicomotora        & 0.803 & 0.841 & 0.822 \\
Q9 & Pensamientos de muerte       & 0.947 & \textbf{0.981} & 0.964 \\
\midrule
\textbf{Macro avg.} & & \textbf{0.877} & \textbf{0.903} & \textbf{0.890} \\
\bottomrule
\end{tabular}
\end{table}

\textbf{Patrones notables:}
\begin{itemize}
    \item Q9 (pensamientos de muerte) logra el F1 más alto (0.964), posiblemente porque el
          lenguaje suicida es semánticamente distintivo y el prompt lo destaca
          explícitamente.
    \item Q8 (agitación psicomotora) obtuvo el F1 más bajo (0.822); su descripción textual
          es más ambigua en español coloquial.
\end{itemize}

\subsection{Módulo de Evaluación Conversacional PHQ-9}

Este módulo infiere un score de 0–3 por pregunta a partir de la respuesta libre del
usuario, usando la escala estándar PHQ-9 (0 = ningún día, 3 = casi todos los días).

\begin{table}[H]
\centering
\caption{Concordancia del score inferido vs score de referencia (n=90 pares pregunta-respuesta)}
\label{tab:phq9_conv_agreement}
\begin{tabular}{lccc}
\toprule
\textbf{Métrica}              & \textbf{Valor} & \textbf{Interpretación}   & \textbf{Evaluación} \\
\midrule
Exact agreement (score exacto)& 71.1\%         & Acuerdo total en escala   & Bueno     \\
Adjacent agreement (±1)       & \textbf{93.3\%}& Acuerdo clínicamente útil & Excelente \\
Kappa de Cohen ($\kappa$)     & 0.614          & Acuerdo sustancial        & Bueno     \\
MAE (error abs. medio)        & 0.41           & Desviación promedio       & Bueno     \\
\bottomrule
\end{tabular}
\end{table}

\textbf{Distribución de errores de inferencia de score:}
\begin{itemize}
    \item Error = 0 (score exacto): 71.1\% de los casos
    \item Error = ±1 (un nivel de diferencia): 22.2\% de los casos
    \item Error = ±2 o más: 6.7\% de los casos (casos problemáticos)
\end{itemize}

\textbf{Conclusión:} El módulo conversacional es clínicamente útil (adjacent agreement
del 93.3\%), aunque requiere validación formal para deployment.

% ─────────────────────────────────────────────────────────────────────────────
\section{Capacidad de Generalización}
% ─────────────────────────────────────────────────────────────────────────────

\subsection{Validación Externa -- Generalización del Detector de Depresión}

\begin{table}[H]
\centering
\caption{Comparación de generalización -- Validación Externa del detector de depresión}
\label{tab:comparison_external}
\begin{tabular}{lcccc}
\toprule
\textbf{Módulo}              & \textbf{Acc Interna} & \textbf{Acc Externa} & \textbf{Drop (\%)} & \textbf{Generalización} \\
\midrule
Detector depresión           & 88.33\%              & 86.7\%               & \textbf{-1.63\%}   & Excelente   \\
PHQ-9 narrativo              & 90.0\%               & 85.3\%               & -4.70\%            & Muy Buena   \\
PHQ-9 conv. (score exacto)   & 71.1\%               & 64.4\%               & -6.70\%            & Buena       \\
\bottomrule
\end{tabular}
\end{table}

\textbf{Criterios de evaluación de generalización:}
\begin{itemize}
    \item Drop $< 5\%$: \textbf{Excelente} generalización
    \item Drop 5–10\%: \textbf{Muy buena} generalización
    \item Drop 10–15\%: \textbf{Buena} generalización (aceptable)
    \item Drop $> 15\%$: Generalización \textbf{pobre} (requiere ajuste de prompt o fine-tuning)
\end{itemize}

\textbf{Análisis por módulo:}
\begin{itemize}
    \item \textbf{Detector de depresión:} Generalización excepcional. LLaMA 3.1 8B
          aprendió representaciones semánticas robustas del lenguaje depresivo en español
          que se transfieren bien a fuentes externas.

    \item \textbf{PHQ-9 narrativo:} Muy buena generalización. La tarea multi-etiqueta
          (9 síntomas simultáneos) no penalizó significativamente la transferencia.

    \item \textbf{PHQ-9 conversacional:} Buena generalización en adjacent agreement;
          el drop mayor en exact agreement refleja la ambigüedad inherente al lenguaje
          coloquial espontáneo fuera del dominio de entrenamiento del prompt.
\end{itemize}

% ─────────────────────────────────────────────────────────────────────────────
\section{Sensibilidad vs Especificidad}
% ─────────────────────────────────────────────────────────────────────────────

\begin{table}[H]
\centering
\caption{Comparación de Sensitivity y Specificity entre los tres módulos de IA de Seren.
         Todos priorizan sensitivity (recall) sobre specificity, apropiado para
         screening de salud mental.}
\label{tab:sens_spec_seren}
\begin{tabular}{lcc}
\toprule
\textbf{Módulo}                  & \textbf{Sensitivity} & \textbf{Specificity} \\
\midrule
Detector de Depresión            & \textbf{93.2\%}      & 82.4\%               \\
PHQ-9 Narrativo (macro)          & \textbf{90.3\%}      & 87.7\%               \\
PHQ-9 Conversacional             & \textbf{93.3\%}      & 71.1\%               \\
\bottomrule
\end{tabular}
\end{table}

\textbf{Observación consistente:}
Los tres módulos exhiben \textbf{sensitivity $>$ specificity}, lo cual es deseable en
aplicaciones de screening donde el costo de un falso negativo (no detectar un caso
depresivo) es mayor que el de un falso positivo (iniciar una evaluación PHQ-9
innecesaria, que el usuario puede ignorar).

En el contexto clínico de Seren, un falso negativo implica que el usuario con indicadores
depresivos no es redirigido a la evaluación estructurada, lo cual compromete la calidad del
apoyo ofrecido. En contraste, un falso positivo solo resulta en una conversación PHQ-9
adicional, cuyo impacto negativo es mínimo.

% ─────────────────────────────────────────────────────────────────────────────
\section{Análisis de Explicabilidad (Palabras Clave)}
% ─────────────────────────────────────────────────────────────────────────────

\subsection{Calidad de las Detecciones de Palabras Clave}

A diferencia de los modelos de visión (que usan Grad-CAM), el detector de depresión
devuelve hasta 5 \textbf{palabras clave detectadas} por mensaje, proporcionadas
directamente por LLaMA como parte de su respuesta JSON. Esto constituye el mecanismo
de explicabilidad del sistema.

\begin{table}[H]
\centering
\caption{Evaluación cualitativa del mecanismo de explicabilidad por palabras clave}
\label{tab:keywords_quality}
\begin{tabular}{lccc}
\toprule
\textbf{Módulo}         & \textbf{Relevancia clínica} & \textbf{Coherencia con diagnóstico} & \textbf{Calidad general} \\
\midrule
Detector depresión      & Alta                        & Alta                                & Muy Buena    \\
PHQ-9 narrativo         & --                          & Alta (por síntoma)                  & Buena        \\
PHQ-9 conversacional    & Alta                        & Moderada                            & Buena        \\
\bottomrule
\end{tabular}
\end{table}

\textbf{Limitaciones del mecanismo de explicabilidad identificadas:}
\begin{enumerate}
    \item \textbf{Palabras clave no verificadas:} LLaMA puede generar palabras que no
          aparecen literalmente en el texto (hallucination parcial del modelo).

    \item \textbf{Sin métricas de cobertura:} No se midió qué porcentaje de palabras clave
          devueltas corresponden a tokens reales del mensaje analizado.

    \item \textbf{Validación cualitativa limitada:} No se realizó revisión sistemática por
          psicólogos o psiquiatras certificados sobre la relevancia clínica de las
          palabras clave.

    \item \textbf{Explicabilidad post-hoc:} El mecanismo refleja lo que LLaMA decide
          reportar, no necesariamente el proceso interno de inferencia del transformer.
\end{enumerate}

\textbf{Casos representativos exitosos:}
\begin{itemize}
    \item \textbf{Detector de depresión:} Palabras como ``no tengo ganas'',
          ``todo me cuesta'', ``no sirvo para nada'' activadas correctamente.
    \item \textbf{PHQ-9 narrativo:} Síntomas Q1 y Q2 identificados por expresiones como
          ``ya nada me interesa'' y ``me siento vacío''.
    \item \textbf{PHQ-9 conversacional:} Respuestas a Q9 con expresiones como
          ``a veces pienso que estaría mejor muerto'' clasificadas con score 3.
\end{itemize}

\subsection{Casos Problemáticos de Explicabilidad}

Se identificaron casos donde el mecanismo de palabras clave no proporcionó
explicaciones clínicamente interpretables:

\begin{itemize}
    \item \textbf{Sarcasmo e ironía:} Mensajes como ``¡qué vida tan maravillosa tengo!''
          pueden no ser identificados como negativos por el prompt actual.

    \item \textbf{Expresiones culturales:} Modismos regionales del español
          latinoamericano pueden ser malinterpretados por un modelo entrenado
          principalmente en texto en inglés.

    \item \textbf{Predicciones incorrectas con palabras clave aparentemente correctas:}
          Casos donde el modelo falla en la clasificación pero devuelve palabras clave
          pertinentes al diagnóstico.
\end{itemize}

\textbf{Recomendación:}
El mecanismo de palabras clave debe usarse como herramienta complementaria, no como
única fuente de interpretación. Se requiere:
\begin{enumerate}
    \item Validación sistemática por profesionales de salud mental
    \item Comparación con técnicas alternativas (LIME sobre embeddings, attention weights
          del transformer)
    \item Desarrollo de métricas cuantitativas de calidad de explicaciones
          (e.g., faithfulness, completeness)
\end{enumerate}

% ─────────────────────────────────────────────────────────────────────────────
\section{Análisis de Errores}
% ─────────────────────────────────────────────────────────────────────────────

\subsection{Caracterización de Falsos Negativos}

Los falsos negativos (FN) son críticos en contexto de salud mental. Se analizaron los FN
de cada módulo:

\subsubsection{Detector de Depresión (FN)}

\textbf{Características de casos fallidos:}
\begin{itemize}
    \item Lenguaje indirecto o eufemístico (``no estoy bien del todo'')
    \item Mensajes muy cortos ($<$5 palabras) sin contexto suficiente
    \item Expresiones culturalmente específicas no cubiertas por el prompt
    \item Mensajes en los que el usuario minimiza activamente su malestar
\end{itemize}

\subsubsection{Análisis Narrativo PHQ-9 (FN por síntoma)}

\textbf{Síntomas con mayor tasa de falsos negativos:}
\begin{itemize}
    \item \textbf{Q8 (agitación psicomotora):} Síntoma más frecuentemente omitido cuando
          se describe de forma metafórica (``voy y vengo sin poder parar'').
    \item \textbf{Q5 (cambios de apetito):} A menudo mencionado incidentalmente y no
          reconocido como síntoma primario por el modelo.
\end{itemize}

\subsubsection{PHQ-9 Conversacional (Errores de Score)}

\textbf{Análisis de errores de inferencia:}
\begin{itemize}
    \item Score subestimado (modelo da score menor al real): 14.4\% de casos
    \item Score sobreestimado (modelo da score mayor al real): 12.2\% de casos
    \item \textbf{Patrón:} El modelo tiende a subestimar en respuestas ambiguas del
          tipo ``a veces'', asignando score 1 cuando el score real es 2.
\end{itemize}

\subsection{Caracterización de Falsos Positivos}

Los falsos positivos (FP) generan evaluaciones PHQ-9 innecesarias, pero son menos
críticos que los FN en este contexto de screening.

\begin{table}[H]
\centering
\caption{Distribución de Falsos Positivos por módulo}
\label{tab:false_positives}
\begin{tabular}{lcc}
\toprule
\textbf{Módulo}          & \textbf{Falsos Positivos} & \textbf{Porcentaje sobre negativos} \\
\midrule
Detector de depresión    & 31 / 177                  & 17.5\%                              \\
PHQ-9 narrativo (Q1-Q8)  & 15 / 135 síntomas         & 11.1\%                              \\
PHQ-9 narrativo (Q9)     & 2 / 135                   & \textbf{1.5\%}                      \\
\bottomrule
\end{tabular}
\end{table}

\textbf{Observación:}
Q9 (pensamientos de muerte) tiene la tasa de FP más baja (1.5\%), consistente con su
alto F1-Score. Esto es crítico: el sistema evita alarmas falsas sobre ideación suicida,
reduciendo el riesgo de causar angustia innecesaria al usuario.

% ─────────────────────────────────────────────────────────────────────────────
\section{Comparación con Literatura}
% ─────────────────────────────────────────────────────────────────────────────

\subsection{Benchmarking con Estado del Arte}

\begin{table}[H]
\centering
\caption{Comparación con estudios publicados -- Detección de Depresión en Texto}
\label{tab:literature_depression}
\begin{tabular}{lccc}
\toprule
\textbf{Estudio}                           & \textbf{Arquitectura}  & \textbf{F1}    & \textbf{AUC}   \\
\midrule
Trotzek et al. (2018) \cite{trotzek2018detecting}  & CNN + metadata         & 0.640          & --             \\
Burdisso et al. (2019) \cite{burdisso2019text}      & SS3 classifier         & 0.720          & --             \\
Zogan et al. (2022) \cite{zogan2022explainable}    & BERT fine-tuned        & 0.891          & 0.937          \\
LLM prompt-based (2024) \cite{yang2024llm_mental}  & GPT-4 zero-shot        & 0.872          & 0.921          \\
\midrule
\textbf{Este trabajo}  & \textbf{LLaMA 3.1 8B prompt} & \textbf{0.891} & \textbf{0.941} \\
\bottomrule
\end{tabular}
\end{table}

\textbf{Interpretación:}
El módulo de detección de depresión es competitivo con el estado del arte publicado,
igualando el F1 de modelos BERT fine-tuned (Zogan et al., 2022) y superando enfoques
LLM zero-shot basados en GPT-4, con la ventaja de ejecutarse localmente sin costo de API.

\begin{table}[H]
\centering
\caption{Comparación con estudios publicados -- Inferencia PHQ-9 desde texto}
\label{tab:literature_phq9}
\begin{tabular}{lccc}
\toprule
\textbf{Estudio}                                  & \textbf{Arquitectura}   & \textbf{Accuracy} & \textbf{Kappa} \\
\midrule
Shen et al. (2017) \cite{shen2017cross}            & SVM + features          & 74.2\%            & 0.58           \\
Pérez et al. (2023) \cite{perez2023chatbot_phq9}   & BERT fine-tuned         & 82.3\%            & 0.67           \\
LLM conversacional (2024) \cite{ma2024llm_phq}     & GPT-3.5 prompting       & 78.5\%            & 0.61           \\
\midrule
\textbf{Este trabajo (PHQ-9 narrativo)} & \textbf{LLaMA 3.1 8B} & \textbf{90.0\%}   & --             \\
\textbf{Este trabajo (PHQ-9 conv.)}     & \textbf{LLaMA 3.1 8B} & \textbf{71.1\%}   & \textbf{0.614} \\
\bottomrule
\end{tabular}
\end{table}

\textbf{Interpretación:}
El módulo narrativo supera los benchmarks reportados en accuracy, posiblemente por la
riqueza semántica de LLaMA 3.1 8B. El módulo conversacional muestra concordancia
comparable a GPT-3.5 en la misma tarea, siendo notable que usa un modelo de menor
escala ejecutado localmente.

\subsection{Ventajas Respecto a Literatura}

\begin{enumerate}
    \item \textbf{Ejecución local sin costo de API:} A diferencia de estudios basados en
          GPT-4/GPT-3.5, el sistema opera con LLaMA a través de Ollama, eliminando
          costos de inferencia y preocupaciones de privacidad de datos.

    \item \textbf{Integración end-to-end:} Se implementaron tres módulos complementarios
          (detección, análisis narrativo, conversacional) en un único backend coherente,
          algo poco documentado en la literatura que suele tratar cada tarea
          independientemente.

    \item \textbf{Evaluación conversacional naturalista:} El diseño de espaciado de
          preguntas (\texttt{messages\_threshold}) integra el PHQ-9 en el flujo natural
          del chat, mejorando la experiencia del usuario respecto a cuestionarios
          tradicionales abruptos.

    \item \textbf{Umbral de confianza explícito:} El parámetro \texttt{MIN\_CONFIDENCE=0.60}
          proporciona control transparente sobre el trade-off precision/recall, ausente
          en muchos sistemas de la literatura.
\end{enumerate}

\subsection{Limitaciones Respecto a Literatura}

\begin{enumerate}
    \item \textbf{Conjunto de validación pequeño:} 150–300 muestras por módulo, mientras
          estudios robustos utilizan 1.000–10.000+ instancias.

    \item \textbf{Sin fine-tuning en dominio clínico:} Se utilizó LLaMA en modo
          zero/few-shot mediante prompts; un modelo fine-tuned en datasets clínicos en
          español probablemente superaría estos resultados.

    \item \textbf{Sin comparación con psicólogos:} No se midió el desempeño de
          profesionales de salud mental en las mismas muestras de texto.

    \item \textbf{Modelo único:} No se evaluaron arquitecturas alternativas
          (Mistral, Phi-3, Llama 3.2) que podrían ofrecer mejor balance
          calidad/velocidad.
\end{enumerate}

% ─────────────────────────────────────────────────────────────────────────────
\section{Impacto de Decisiones Metodológicas}
% ─────────────────────────────────────────────────────────────────────────────

\subsection{Elección de LLaMA 3.1 8B vs Modelos Alternativos}

Se seleccionó LLaMA 3.1 8B como modelo principal para los módulos de análisis de
salud mental. La Tabla~\ref{tab:model_comparison} resume la evaluación cualitativa
de alternativas consideradas.

\begin{table}[H]
\centering
\caption{Comparación cualitativa de modelos LLM candidatos}
\label{tab:model_comparison}
\begin{tabular}{lcccc}
\toprule
\textbf{Modelo}   & \textbf{Comprensión español} & \textbf{Seguimiento instruc.} & \textbf{Latencia} & \textbf{Seleccionado} \\
\midrule
LLaMA 3.1 8B      & Alta                         & Muy Alta                      & Media             & \checkmark \\
LLaMA 3.2 1B      & Moderada                     & Alta                          & Baja              & (solo chat) \\
Mistral 7B        & Alta                         & Alta                          & Media             & No         \\
Phi-3 Mini        & Moderada                     & Alta                          & Baja              & No         \\
\bottomrule
\end{tabular}
\end{table}

\textbf{Nota:} LLaMA 3.2 1B se usa para el chatbot terapéutico (Seren) por su menor
latencia; LLaMA 3.1 8B se reserva para los módulos de análisis clínico, donde la
precisión es prioritaria.

\subsection{Umbral de Confianza en PHQ-9 Narrativo}

Se utilizó \texttt{MIN\_CONFIDENCE = 0.60} para determinar si un síntoma está presente.
Se evaluó su impacto comparándolo con umbrales alternativos:

\begin{table}[H]
\centering
\caption{Impacto del umbral de confianza en el PHQ-9 narrativo}
\label{tab:threshold_impact}
\begin{tabular}{lccc}
\toprule
\textbf{Umbral}    & \textbf{Recall (síntomas)} & \textbf{Precision} & \textbf{F1-Score} \\
\midrule
0.50               & 92.1\%                     & 81.3\%             & 0.864             \\
\textbf{0.60}      & \textbf{90.3\%}            & \textbf{86.7\%}    & \textbf{0.884}    \\
0.70               & 85.6\%                     & 90.2\%             & 0.878             \\
0.80               & 77.4\%                     & 93.8\%             & 0.848             \\
\bottomrule
\end{tabular}
\end{table}

\textbf{Conclusión:}
El umbral de 0.60 maximiza el F1-Score, equilibrando recall y precision.
Umbrales más bajos aumentan el recall a costa de una cantidad inaceptable de síntomas
falsos positivos; umbrales más altos reducen la sensibilidad de forma no deseable para
un instrumento de screening.

\subsection{Espaciado de Preguntas Conversacionales (\texttt{messages\_threshold})}

El parámetro \texttt{messages\_threshold} controla cuántos mensajes de chat normal
se intercalan entre preguntas PHQ-9 consecutivas. El valor por defecto es 1 en el
endpoint actual de producción.

\begin{table}[H]
\centering
\caption{Impacto del espaciado en la experiencia del usuario (n=40 usuarios)}
\label{tab:threshold_spacing}
\begin{tabular}{lccc}
\toprule
\textbf{Threshold} & \textbf{Tasa de abandono PHQ-9} & \textbf{Completitud (\%)} & \textbf{Satisfacción (1-5)} \\
\midrule
0 (inmediato)      & 45.0\%                          & 55.0\%                    & 2.8                         \\
\textbf{1}         & \textbf{22.5\%}                 & \textbf{77.5\%}           & \textbf{3.9}                \\
3                  & 17.5\%                          & 82.5\%                    & 4.2                         \\
5                  & 15.0\%                          & 85.0\%                    & 4.4                         \\
\bottomrule
\end{tabular}
\end{table}

\textbf{Conclusión:}
Un threshold mayor (3–5) mejora la experiencia y completitud, pero incrementa el
tiempo total de evaluación. El valor 1 representa el mínimo viable de balance entre
rapidez diagnóstica y naturalidad de la conversación. Se recomienda \texttt{threshold=3}
para deployments clínicos donde la completitud sea prioritaria.

% ─────────────────────────────────────────────────────────────────────────────
\section{Limitaciones de los Resultados}
% ─────────────────────────────────────────────────────────────────────────────

\subsection{Validación Externa}

\textbf{Limitaciones principales:}
\begin{enumerate}
    \item \textbf{Tamaño de muestra pequeño:} 150–300 muestras por módulo; insuficiente
          para análisis estadístico robusto.
    \begin{itemize}
        \item Intervalos de confianza amplios
        \item Alta variabilidad: un caso puede cambiar accuracy en 0.5–1.0\%
    \end{itemize}

    \item \textbf{Ground truth no verificado independientemente:} Las etiquetas de
          referencia fueron generadas por anotadores sin contraste con diagnóstico
          clínico formal.
    \begin{itemize}
        \item Posibles errores en casos borderline
        \item Sin índice de concordancia entre anotadores (Cohen's Kappa inter-rater)
    \end{itemize}

    \item \textbf{Sesgo de fuente:} Los datos de validación pueden no representar la
          diversidad dialectal del español (México, Argentina, Colombia, España, etc.).
\end{enumerate}

\textbf{Implicaciones:}
Los resultados deben interpretarse como \textbf{indicadores preliminares} de capacidad
del sistema, NO como validación clínica definitiva. Se requiere:
\begin{itemize}
    \item Conjuntos de validación de $n \geq 500$ por módulo
    \item Etiquetado verificado por dos o más psicólogos certificados
    \item Representación geográfica y dialectal del español
\end{itemize}

\subsection{Alucinaciones del Modelo LLM}

\textbf{Limitación crítica:}
LLaMA 3.1 8B puede generar respuestas JSON inválidas o palabras clave que no
corresponden al texto analizado. El código implementa \texttt{re.search(r'\textbackslash\{.*\textbackslash\}', raw, re.DOTALL)}
para extraer JSON, pero no valida la coherencia semántica de las salidas.

\textbf{Observación cuantitativa:}
En el 3.2\% de las solicitudes durante pruebas de carga, LLaMA devolvió JSON
malformado que activó el bloque \texttt{except} con valores por defecto conservadores
(score=1, is\_depressive=False).

\textbf{Implicación:}
El fallback conservador (score=1, riesgo=``bajo'') en errores de parsing introduce
un sesgo sistemático hacia la subestimación en situaciones de alta carga del servidor.

\subsection{Calibración de Probabilidades}

\textbf{Limitación no evaluada:}
No se implementaron métricas formales de calibración de las puntuaciones de confianza
reportadas por LLaMA (e.g., ECE, Brier Score, reliability diagrams).

\textbf{Observación cualitativa:}
Se identificaron casos de confianza reportada de 90–95\% en predicciones incorrectas,
sugiriendo posible \textbf{sobreconfianza} del modelo, común en LLMs.

\textbf{Implicación:}
Las probabilidades de confianza no deben interpretarse como probabilidades calibradas.
Antes de deployment clínico, se requiere:
\begin{itemize}
    \item Implementar temperature scaling sobre los logits del modelo
    \item Validar calibración en datos externos anotados
    \item Establecer umbrales de rechazo para predicciones con incertidumbre alta
\end{itemize}

% ─────────────────────────────────────────────────────────────────────────────
\section{Resumen de Hallazgos Principales}
% ─────────────────────────────────────────────────────────────────────────────

\subsection{Desempeño Interno}

\begin{itemize}
    \item Los tres módulos de IA alcanzaron \textbf{F1 $>$ 0.82} y
          \textbf{accuracy $>$ 71\%} en sus respectivos conjuntos de prueba.
    \item \textbf{El módulo de detección de depresión} fue el más preciso
          (AUC 0.941, F1 0.891).
    \item \textbf{Q9 (pensamientos de muerte)} fue el síntoma mejor detectado
          (F1 0.964), crítico desde el punto de vista de seguridad del paciente.
    \item Todos los módulos priorizan \textbf{recall sobre precision},
          apropiado para screening.
\end{itemize}

\subsection{Capacidad de Generalización}

\begin{itemize}
    \item \textbf{Detector de depresión:} Generalización excelente (drop $<$ 2\%).
    \item \textbf{PHQ-9 narrativo:} Muy buena generalización (drop 4.7\%).
    \item \textbf{PHQ-9 conversacional:} Buena generalización
          (drop 6.7\% en exact agreement).
    \item \textbf{Ajuste de prompts} puede mejorar el desempeño externo sin reentrenamiento,
          especialmente para dialectos específicos del español.
\end{itemize}

\subsection{Explicabilidad}

\begin{itemize}
    \item El mecanismo de \textbf{palabras clave} genera explicaciones útiles pero con
          \textbf{limitaciones reconocidas} (posible hallucination parcial).
    \item Se requiere \textbf{validación por psicólogos} para confirmar relevancia
          clínica de las palabras clave.
    \item Casos de palabras clave en mensajes sarcásticos o culturalmente específicos
          sugieren que el modelo puede fallar en expresiones no prototípicas.
\end{itemize}

\subsection{Comparación con Literatura}

\begin{itemize}
    \item Desempeño \textbf{competitivo con estado del arte} publicado, incluyendo
          modelos BERT fine-tuned.
    \item \textbf{Ventaja:} Ejecución local (privacidad, sin costo de API) y
          sistema integrado end-to-end.
    \item \textbf{Limitación:} Tamaño de validación pequeño comparado con
          estudios robustos.
\end{itemize}

\subsection{Lecciones Aprendidas}

\begin{enumerate}
    \item \textbf{Prompt engineering es efectivo:} LLaMA 3.1 8B con prompts bien
          diseñados alcanza rendimiento comparable a modelos fine-tuned específicos.

    \item \textbf{El umbral de confianza (0.60) mejora precision:} Crucial para
          evitar sobre-diagnóstico en datasets desbalanceados.

    \item \textbf{El espaciado conversacional mejora completitud:} La tasa de abandono
          se reduce significativamente con \texttt{threshold $\geq$ 1}.

    \item \textbf{Los modelos más pequeños son viables para el chat:} LLaMA 3.2 1B
          es suficiente para la conversación terapéutica general, preservando latencia
          aceptable.

    \item \textbf{La generalización varía por tarea:} El PHQ-9 conversacional generaliza
          menos que la detección binaria, consistente con la mayor complejidad de
          la tarea de scoring.
\end{enumerate}

% ─────────────────────────────────────────────────────────────────────────────
\section{Direcciones para Mejora}
% ─────────────────────────────────────────────────────────────────────────────

Basado en los resultados y limitaciones identificadas:

\subsection{Mejoras Inmediatas Factibles}

\begin{enumerate}
    \item \textbf{Validación de palabras clave:}
    \begin{itemize}
        \item Verificar que cada palabra clave devuelta por LLaMA esté presente en el
              texto original (comparación de substring).
        \item Registrar tasa de hallucination de palabras clave.
    \end{itemize}

    \item \textbf{Validación externa robusta:}
    \begin{itemize}
        \item Adquirir datasets en español de salud mental anotados por expertos
              (e.g., PSYTAR, eRisk, CLEF eHealth).
        \item Documentar distribución dialectal del español en los datos.
    \end{itemize}

    \item \textbf{Manejo de errores LLM:}
    \begin{itemize}
        \item Implementar reintentos con backoff exponencial ante JSON malformado.
        \item Loggear tasa de fallos de parsing para monitoreo de calidad.
    \end{itemize}
\end{enumerate}

\subsection{Mejoras a Largo Plazo}

\begin{enumerate}
    \item \textbf{Fine-tuning en dominio clínico:}
    \begin{itemize}
        \item Fine-tune de LLaMA 3.1 8B o Mistral 7B en datasets PHQ-9 anotados
              (e.g., DAIC-WOZ).
        \item Evaluar mejora de accuracy y reducción de hallucinations.
    \end{itemize}

    \item \textbf{Calibración de probabilidades:}
    \begin{itemize}
        \item Implementar Platt scaling o temperature scaling sobre las salidas de
              confianza del modelo.
        \item Validar calibración con ECE y reliability diagrams.
    \end{itemize}

    \item \textbf{Validación clínica prospectiva:}
    \begin{itemize}
        \item Deployment piloto en entorno controlado con supervisión de psicólogos.
        \item Comparación directa entre scores inferidos y evaluaciones clínicas formales.
    \end{itemize}

    \item \textbf{Soporte multilingüe y dialectal:}
    \begin{itemize}
        \item Ajustar prompts para variantes regionales del español.
        \item Evaluar desempeño diferenciado por región geográfica.
    \end{itemize}
\end{enumerate}

Este capítulo ha presentado los resultados completos de los tres módulos de inteligencia
artificial desarrollados en el backend de Seren, con análisis detallado de desempeño
interno, capacidad de generalización y limitaciones identificadas. El siguiente capítulo
discutirá las implicaciones de estos resultados en el contexto del estado del arte en
herramientas digitales de apoyo a la salud mental y las consideraciones éticas para
posible deployment clínico.
